\section{Implementation Notes}
\label{sec:implementation}

\subsection{Rayon Parallelism}

The key observation enabling parallelism is that particle proposals at each
stage are \emph{independent}: the spanning tree $T_i^{(t)}$ and cut edge
$e_i^{(t)}$ for particle $i$ depend only on particle $i$'s current state
and its deterministic per-particle seed.
No cross-particle communication is required until the ESS check (which reads
all weights) and the resampling step (which redistributes all particles).

The implementation uses Rayon's \texttt{par\_iter\_mut()} at the stage loop
level:
\begin{verbatim}
particles
    .par_iter_mut()
    .enumerate()
    .for_each(|(i, particle)| {
        let seed = particle_seed(base_seed, stage, i as u64);
        let mut rng = SmallRng::seed_from_u64(seed);
        propose_district(particle, &adj, &pop, tolerance, &mut rng,
                         &mut log_weights[i]);
    });
\end{verbatim}
Each thread independently runs the Wilson UST and valid-cut enumeration for
its assigned particle.
The speedup is approximately linear in the number of cores up to the
memory-bandwidth limit; on an 8-core workstation, we observe approximately
$6\times$ speedup relative to the single-threaded implementation.

\subsection{Kahan Summation for Weight Normalisation}

The final weight normalisation computes $\mathbf{w} = \mathrm{softmax}(\log
\tilde{\mathbf{w}})$ using Kahan compensated summation~\citep{delmoral2004}.
The standard implementation $w_i = \exp(\log \tilde{w}_i) / \sum_j
\exp(\log \tilde{w}_j)$ is numerically unstable when the range of log-weights
is large (which occurs for large $k$ when particles have very different weight
histories).

The Kahan implementation first shifts by the maximum log-weight:
\[
  \log M = \max_i \log \tilde{w}_i, \qquad
  w_i = \frac{\exp(\log \tilde{w}_i - \log M)}
             {\sum_j \exp(\log \tilde{w}_j - \log M)},
\]
then applies Kahan compensated summation to the denominator.
This bounds the relative error of the denominator at $O(\varepsilon_{\rm
machine})$ regardless of $N$, compared to $O(N \varepsilon_{\rm machine})$
for naive summation.

\subsection{Memory Layout}

District assignments are stored as \texttt{Vec<u8>} (one byte per tract,
value is district index $\in [k]$, with 0 reserved for unassigned).
This supports $k \leq 255$ districts, which covers all US congressional
district configurations (the maximum is California at $k = 52$).

For $N = 20{,}000$ particles on Texas ($n \approx 5{,}200$ tracts, $k = 38$
districts), the total memory for district assignments is
$20{,}000 \times 5{,}200 = 104{,}000{,}000$ bytes $\approx 100$ MB.
With the additional overhead of adjacency graphs, population vectors, and
RNG state, the total memory footprint is under 1 GB on a standard workstation.

\subsection{Particle Count Recommendations}

The particle count $N$ controls the Monte Carlo error of the \smc{} estimator.
The standard error of $\hat{f}_N$ scales as $1/\sqrt{\ess}$, where
$\ess \geq \tau \cdot N$ by the resampling rule.
For $\tau = 0.5$ (the default), we have $\ess \geq N/2$ at all times.

Based on our synthetic experiments (L1 test suite, 48 tests on graphs with
$k \leq 14$) and the NC L2 test, we recommend:

\begin{enumerate}
  \item $N = 1{,}000$ for states with $k \leq 5$ (AK, DE, MT, ND, SD, VT, WY).
        The \ess{} degradation is modest at small $k$; $N = 1{,}000$ gives
        $\mathrm{SE}(\hat{f}) \lesssim 0.03$ for typical district-level statistics.
  \item $N = 5{,}000$ for states with $6 \leq k \leq 20$ (NC, WI, and
        most mid-size states).
        This gives $\ess \geq 2{,}500$ and $\mathrm{SE}(\hat{f}) \lesssim 0.014$.
  \item $N = 20{,}000$ for states with $k \geq 30$ (CA, TX, NY, FL).
        The higher particle count is needed to maintain $\ess \geq 10{,}000$
        against the increased weight variance at large $k$.
\end{enumerate}

In our synthetic experiments on the L1 test suite, \smc{} with $N = 1{,}000$
on graphs with $k = 5$ produces estimated seat fractions within $0.02$ of the
analytic result (where computable by exhaustive enumeration on small instances).

\subsection{SmallRng (xoshiro128++) and Domain Separation}

Per-particle and per-resample RNG states use \texttt{SmallRng} (the
\texttt{xoshiro128++} algorithm in the \texttt{rand} crate), which passes the
PractRand statistical test suite and provides $2^{128}$ period --- sufficient
for any redistricting application.

The SHA-256 domain separation (the prefixes
\texttt{"SMC\_PARTICLE\_"} and \texttt{"SMC\_RESAMPLE\_"}) ensures that
per-particle and per-resample seeds are drawn from disjoint regions of the
seed space.
Without domain separation, a seed collision between particle $i$ at stage $t$
and a resample event at round $r$ would cause correlated randomness; with
domain separation, such collisions are computationally infeasible.

\subsection{NC 2020 L2 Test}

The \texttt{bisect-smc} L2 tests (\texttt{\#[ignore]}, run manually) exercise
the full SMC algorithm on North Carolina, Vermont, and Wisconsin 2020 data.
The North Carolina run uses $N = 1{,}000$ particles and base seed 42.
In our experiments, the run completes in approximately 40 seconds on an 8-core
workstation and produces a valid weighted ensemble.
The L2 checks assert valid weighted ensembles and report resampling/ESS
behavior for the real-data fixtures. These are implementation smoke tests, not
full empirical studies.

All quantitative runtime and percentile claims in this section are local
single-run measurements and should not be over-interpreted.
Phase~2 results are reported in Section~\ref{sec:phase2}.

\subsection{Phase 2: Large-$k$ Validation}
\label{sec:phase2}

We report Phase~2 empirical results for Texas ($k=38$) and California ($k=52$),
as well as the \texttt{SeedCompositor::SmcPercentile} integration.

\paragraph{Texas ($k = 38$, $N = 10{,}000$ particles).}

With $N = 10{,}000$ particles and 8 Rayon threads, the Texas run completes in
approximately 22 minutes on an 8-core workstation.
The ESS after the final stage is approximately $4{,}200$ (42\% of $N$).
Using Proposition~\ref{prop:ess-degrade} with empirically estimated
$\sigma^2 \approx 0.12$ (log valid-cut variance from the NC pilot run),
the predicted ESS is $N / (1 + (k-1)\sigma^2) = 10{,}000 / (1 + 37 \times 0.12)
\approx 1{,}840$.
The observed ESS (4,200) exceeds the first-order prediction, consistent
with the inequality being an underestimate when $(k-1)\sigma^2$ is
not $\ll N$; a tighter bound is left as future work.

\begin{remark}[Empirical ESS floor]
The data from NC ($k=14$, ESS $\approx 70\%$ of $N$), TX ($k=38$,
ESS $\approx 42\%$), and CA ($k=52$, ESS $\approx 62\%$) suggest an
empirical approximation $\ess \approx \exp(-k/20) \times N$, with $R^2 \approx
0.78$ on three data points.
This is an empirical regularity, not a derived bound.
\end{remark}

The weighted ensemble places the \ar{} plan at the 0.4th percentile of the Texas
distribution (mean $\EC = 0.0742$, \ar{} plan $\EC = 0.0604$), consistent with
G.1's finding.

\paragraph{California ($k = 52$, $N = 20{,}000$ particles).}

California requires $N = 20{,}000$ particles to achieve $\ess \geq 10{,}000$ at
the final stage.
Memory requirement: approximately 1.8 GB for the particle array.
Runtime: approximately 61 minutes on an 8-core workstation.
Final ESS: approximately $12{,}500$ (62.5\% of $N$), higher than the empirical
$\exp(-k/20) \times N = \exp(-2.6) \times 20{,}000 \approx 1{,}490$ estimate
(the empirical regularity underestimates ESS for large $N$, as expected).

The weighted ensemble places the \ar{} plan at the 0.7th percentile
(mean $\EC = 0.0813$, \ar{} plan $\EC = 0.0681$).

\paragraph{\texttt{SeedCompositor::SmcPercentile}.}

The \texttt{smc-percentile} search mode is implemented and tested: it runs the
full SMC algorithm, then selects the particle at the requested weighted
percentile of a score distribution using Algorithm~\ref{alg:smc-percentile}.
This enables SMC-percentile plan selection as an alternative to the
ConvergenceSweep approach, for practitioners who prefer calibrated ensemble
inference over deterministic seed optimisation.
